\section{Threats to validity}
\label{sec:threats}

\noindent\textbf{Data and scope.} The collections represent local
demonstration and benchmark workloads. Developers authored or expanded the
questions and gold annotations. The routing benchmark includes 52 expanded
questions, with 47 questions concentrated in three English domains.
Same-question reruns test repeatability under recorded changes to the
environment: across the three campaigns, the A0h containment count moved
from 52 to 51 to 48 and the A0 count from 37 to 39 to 36, while the
paired A0h-over-A0 direction held in each record. The campaigns are
reported separately and are not pooled. An independent question set would
test transfer. The storage
reconciliation examined 33 of 322 parts, and the four out-of-domain probes
characterise individual responses.

\noindent\textbf{Measurement and comparisons.} The content metric counts a
literal target string anywhere in the pack, including in a sibling
document. In a post-hoc split, the complete nonempty target occurs in 29
questions after lowercasing both strings. In the Main comparison, A0 and
A0h matched 16/29 and 21/29 in this subset, and 20/43 and 27/43 in the
remainder. These are descriptive, non-random subsets, not a causal test of
overlap; lack of an exact match does not rule out other lexical cues.
Gold-ID checks search answers and recorded trace text; the saved grid
questions do not contain their corresponding IDs. Nevertheless, ID presence
does not establish source resolution or semantic correctness, and
claim-to-source support requires its own evaluation; citation markers only
indicate recognised formatting. The full-read comparator is a mechanical
scan of the corpus. Content arms change several components, including candidate
construction and reading scope. A0, A1b and A3 can use the gold domain when
resolving references. A3 returned candidates from outside the evaluated
libraries. The traces do not quantify the description fallback or locate
all target matches by source. These contrasts apply to the tested
configurations, not to navigation alone. The two ten-question answer grids use one
model and one run per arm, with the common-substrate grid controlling which
libraries each interface uses.

\noindent\textbf{Statistical and operational validity.} The paired
comparisons are exploratory. Wilson intervals summarise uncertainty
conditional on the question set and observed run. Candidate- and span-level
tests can understate uncertainty by treating within-query observations as
independent. Some historical reference-mode inputs, judge inputs and code
hashes are missing. The verifier results use a historical implementation
and support replay of the released records. The lifecycle assertions and
concurrency sweep cover one installation. The experience approval path,
multi-host reliability, and crash recovery remain open.
